
% === Pali ===

\cleartoverso

\vspace*{30mm}

\begin{verse}

\verseref{1}
kāmaṁ kāmayamānassa\\
tassa ce taṁ samijjhati\\
addhā pītimano hoti\\
laddhā macco yadicchati

\verseref{2}
assa ce kāmayānassa\\
chandajātassa jantuno\\
te kāmā parihāyanti\\
sallaviddhova ruppati

\verseref{3}
yo kāme parivajjeti\\
sappasseva padā siro\\
somaṁ visattikaṁ loke\\
sato samativattati

\verseref{4}
khettaṁ vatthuṁ hiraññaṁ vā\\
gavassaṁ dāsaporisaṁ\\
thiyo bandhū puthu kāme\\
yo naro anugijjhati

\end{verse}

% === Slovenian ===

\chapter[Kāma Sutta]{{kama-sutta-gray.png}{ucenje-2.png}}

\begin{verse}

\vFirst
\includegraphics[scale=0.3]{ce-smrtniku-gray.png}

% \verseref{1}
% Če smrtniku uspe\\
% uresničiti svoje čutne želje,\\
% postane vesel v svojem srcu,\\
% saj je dobil tisto, kar si je želel.

\verseref{2}
Če za tega človeka\\
s čutnostjo iz želja rojeno,\\
ti predmeti čutnosti oslabijo,\\
je v taki stiski, kot da je preboden s trnom.

\verseref{3}
Kdor pa se čutnosti izogiba\\
kot stopalo kačji glavi,\\
tak človek, ki je pozoren,\\
preseže posvetno navezanost.

\verseref{4}
Tistega, ki je pohlepen\\
na zemljo, lastnino, zlato,\\
krave, konje, sužnje, služabnike, ženske,\\
odnose in raznovrstne čutnosti,

\end{verse}

% === Pali ===

\clearpage
\begin{verse}

\verseref{5}
abalā naṁ balīyanti\\
maddantenaṁ parissayā\\
tato naṁ dukkhamanveti\\
nāvaṁ bhinnamivodakaṁ

\verseref{6}
tasmā jantu sadā sato\\
kāmāni parivajjaye\\
te pahāya tare oghaṁ\\
nāvaṁ sitvāva pāragūti

\end{verse}

% === Slovenian ===

\clearpage
\begin{verse}

\verseref{5}
premagajo nemočni,\\
težave ga potisnejo k tlom,\\
tako, da neprijetnost pride k njemu,\\
kot voda v polomljen čoln.

\verseref{6}
Zato naj se tisti, ki je vedno pozoren, izogne čutnosti;\\
ko jo bo opustil, bo prečkal to poplavo,\\
tako kot tisti, ki je s črpanjem vode iz svojega čolna,\\
uspel priti na drugi breg.

\end{verse}

